% =====================================================================
% PATCHES PARA O PAPER — substituições prontas a colar.
% Cada bloco indica ONDE substitui e PORQUÊ.
% Base: experiência 16B, 27 projetos / 486 módulos.
% =====================================================================


% ---------------------------------------------------------------------
% [1] ABSTRACT — último parágrafo
% PORQUÊ: a versão atual promete "significantly higher yield of valid,
% executable tests" (na verdade é empate: 11.255 vs 11.235) e "maximizes
% statement coverage ... " sem ressalva (o Pynguin ganha na média por
% projeto). O que os dados sustentam é mais forte e é verdadeiro.
% ---------------------------------------------------------------------
We evaluate MARTA on 27 open-source Python projects (486 modules) against a
search-based generator and a state-of-the-art single-prompt LLM baseline running on
the identical semantic context. Across the benchmark's modules, MARTA achieves the
highest line+branch coverage and the highest mutation score of the three systems,
while its generation phase consumes half the wall-clock time of the single-prompt
pipeline. The gain comes less from producing more tests than from producing tests that
actually verify:
only 0.8\% of MARTA's tests contain no assertion at all, against 7.1\% for the
single-prompt baseline and 38.3\% for the search-based generator, whose structural
coverage is therefore accrued in substantial part by tests that execute code without
checking it. We conclude that decoupling planning from syntax yields suites that
detect more faults per unit of coverage, and that the benefit of such decomposition is
better assessed by fault detection than by structural coverage alone.


% ---------------------------------------------------------------------
% [2] INTRODUÇÃO — 4.º item da lista de contribuições
% PORQUÊ: "substantially higher yields of valid, executable tests,
% statement and branch coverage, and mutation score than single-prompt AND
% search-based baselines" é falso para cobertura vs Pynguin por projeto.
% ---------------------------------------------------------------------
\item An empirical study on a state-of-the-art Python benchmark of 27 projects
(486 modules) showing that the decoupled pipeline produces substantially stronger test
oracles and achieves the highest coverage and mutation score among the evaluated
systems, while consuming roughly half the tokens and wall-clock time of the
single-prompt architecture. We further analyse where each paradigm is most effective,
relating the residual differences to the capacity of the underlying model rather than
to the generation strategy alone.


% ---------------------------------------------------------------------
% [3] §4.3 IMPLEMENTATION DETAILS — parágrafo dos modelos e temperatura
% PORQUÊ: (a) o 236B foi substituído pelo Qwen2.5-Coder-32B por
% restrições de quota (o 236B exige 4 GPUs, ou seja 4x o consumo);
% (b) o paper afirma T=0.2 para ambos, mas o baseline corre a T=0.0
% (o seu valor por omissão) — o artefacto é público e contradiz.
% SUBSTITUI os dois parágrafos que começam em "The pipeline is evaluated
% at two model scales..." e "Following established practices...".
% ---------------------------------------------------------------------
The pipeline is evaluated with two models to assess how the architectural benefits
transfer across model capacity. The primary evaluation, covering all 27 projects, uses
\textbf{DeepSeek-Coder-V2-Lite (16B)}, a mixture-of-experts model that activates only
2.4B parameters per token. To probe whether the observed benefits are sensitive to
model capacity, we additionally run both MARTA and the single-prompt baseline with
\textbf{Qwen2.5-Coder-32B}, a dense model that applies its full 32B parameters to every
token and is therefore a substantially stronger reasoner per unit of inference.

This second configuration is deliberately exploratory. Serving a dense 32B model
locally proved roughly an order of magnitude slower per project than the 16B
mixture-of-experts model, and a full 27-project run for both systems would have
exceeded our entire remaining allocation on the shared cluster. We therefore restrict
it to a subset of ten projects (31 modules), fixed before observing any 32B result and
deliberately balanced: it contains projects on which MARTA leads the baseline at 16B as
well as projects on which the baseline leads, so that a change in the gap cannot be an
artefact of the selection. We treat this configuration as directional evidence about
the role of model capacity rather than as a second full evaluation, and we draw no
scaling conclusions from two points.

Both models are served locally via Ollama, and the retrieval-augmented context module
uses the \texttt{BAAI/bge-large-en-v1.5} sentence embedder across all configurations.
All experiments were conducted on the Deucalion HPC cluster, on a single x86/GPU node
with one NVIDIA A100 (40\,GB) GPU, 32 CPU cores, and 200\,GB of RAM. Because all
inference is served locally rather than through a hosted API, the runtimes we report
reflect genuine on-device inference cost (model loading, prompt evaluation, and token
decoding) rather than network latency or third-party rate limits. Since Phase~1 is
cached and shared by both LLM-based systems, we report the cost of the generation phase
separately from that of context building, and supplement wall-clock time with token
consumption so that the comparison does not depend on the serving stack.

Regarding decoding, each tool is run at the temperature configured by its authors:
$T=0.2$ for MARTA and $T=0.0$ for the single-prompt baseline. Values at or near zero
maximize reproducibility but are known to trap models in degenerate, repetitive
loops \cite{holtzman2020curious}, which is particularly detrimental inside a ReAct
self-healing loop that requires the model to propose \emph{alternative} corrections
rather than re-emitting identical invalid syntax; higher temperatures, conversely,
increase the risk of hallucinated identifiers. We report this asymmetry explicitly
and revisit it in Section~6.1, as it is the one configuration parameter that is not
held constant across the two LLM-based systems. Both systems share an identical
repair budget of $N=3$ LLM generation attempts per test file (one initial generation
followed by at most two repairs), the same Phase~1 semantic context, and the same
execution environment.


% ---------------------------------------------------------------------
% [4] §6.1 THREATS — INTERNAL VALIDITY (substituição integral)
% PORQUÊ: o texto atual afirma "every configuration is repeated across
% multiple independent runs" — foi feita UMA run por ferramenta. E omite
% a assimetria de temperatura.
% ---------------------------------------------------------------------
\subsection{Internal Validity}
Internal validity concerns whether the observed improvements can be attributed to our
architectural treatment rather than to confounding factors. We designed the
evaluation so that the decoupled architecture is, as far as possible, the only
independent variable: the single-prompt baseline and MARTA share the identical
Phase~1 semantic context (the Rich Type Context), the same underlying LLM, the same
repair budget ($N=3$), and the same execution environment. One parameter is not held
constant: each system runs at the decoding temperature configured by its authors
($T=0.2$ for MARTA, $T=0.0$ for the baseline). We report this explicitly rather than
retrofitting a common value, since altering the baseline's configuration would
misrepresent the system as published; the direction of its effect is not obvious, as
a lower temperature favours reproducibility but is also associated with repetitive
outputs inside repair loops.

A second threat concerns sampling variance. Because local inference for the full
benchmark is computationally expensive, each configuration was executed once rather
than repeated; we therefore report per-project distributions (mean and median) and
assess differences with the Wilcoxon signed-rank test over the paired per-project
results, but we cannot report confidence intervals over repeated runs, and
project-level outliers should be interpreted with corresponding caution.

To prevent instrumentation bias, all metrics are computed with established libraries
(\texttt{coverage.py} for structural adequacy and \texttt{mutmut} for mutation
analysis) and, critically, through a single measurement pipeline applied uniformly to
all three tools rather than through each tool's own reporting. This proved necessary:
measuring coverage with each project's own \texttt{coverage.py} configuration would
have silently excluded entire target modules in one subject, and executing an entire
generated suite in a single \texttt{pytest} invocation allows one pathological test to
abort collection and collapse the measured coverage of a whole project. We therefore
execute test suites in batches and with a uniform coverage configuration.

A further threat is prompt-engineering effort, as the agent prompts may have received
more refinement than the baseline; we mitigate this by reusing the baseline's original
generation prompt from TEST4PY \cite{test4py2025} unchanged. Finally, because the
benchmark consists of public open-source projects, partial memorization by the LLM
cannot be excluded; as both systems use the same model, this would affect absolute
scores rather than the relative comparison on which our conclusions rest.


% ---------------------------------------------------------------------
% [5] §6.2 THREATS — EXTERNAL VALIDITY (parágrafo dos modelos)
% PORQUÊ: 236B -> Qwen-32B; e a conclusão sobre capacidade do modelo é
% baseada em DOIS pontos, o que permite afirmar direção, não escala.
% ---------------------------------------------------------------------
Regarding the model, our experiments use two open-weights models
(DeepSeek-Coder-V2-Lite 16B and Qwen2.5-Coder-32B), selected to differ in effective
reasoning capacity while remaining servable on a single GPU. Two points are
sufficient to establish a direction but not a scaling law: our observations about how
the architectural benefit varies with model capacity should be read as indicative
rather than as a characterized trend. The magnitude of the decoupling benefit may
also differ for other model families or for frontier API models, which we did not
evaluate.


% ---------------------------------------------------------------------
% [6] §6.3 THREATS — CONSTRUCT VALIDITY (parágrafo adicional no fim)
% PORQUÊ: documentar as limitações de medição que encontrámos.
% ---------------------------------------------------------------------
Two measurement limitations warrant explicit mention. First, mutation analysis
attributes each mutant to the module that contains it and executes it against the
tests targeting that module; a mutant that would only be detected by a test of a
different module is therefore counted as surviving. We consider this the more
faithful attribution, but it yields systematically lower scores than executing the
entire suite against every mutant. Second, our weak-assertion criterion is
deliberately conservative: it flags only assertions that cannot discriminate
behaviour (constants, type and \texttt{isinstance} checks, and assertions over mock
internals) and does \emph{not} penalize comparisons against literal values, since
these are the canonical form of a correct oracle. The reported weak-assertion rates
are consequently lower bounds.


% ---------------------------------------------------------------------
% [7] §4.3 — LISTA DE BASELINES (substituição integral do parágrafo
%     introdutório e da lista com bullets)
% PORQUÊ: o texto atual anuncia "three distinct state-of-the-art
% paradigms" e inclui o CoverUp, que NÃO foi executado. Chamar baseline
% a um sistema que não se corre é indefensável em revisão. Passam a ser
% DUAS baselines controladas; o CoverUp é reposicionado como referência
% publicada, o que é o seu estatuto real.
% ---------------------------------------------------------------------
We compare MARTA against two baselines, both re-executed by us under identical
conditions, each representing a distinct generation paradigm:
\begin{itemize}
    \item \textbf{Pynguin} \cite{lukasczyk2022pynguin}: the state-of-the-art
    Search-Based Software Testing (SBST) framework for Python, representing
    generation without any language model.
    \item \textbf{Single-Prompt Baseline}: the generation module of TEST4PY
    \cite{test4py2025}, executed on the exact same Phase~1 semantic context, the same
    model, and the same repair budget as MARTA. Because the two differ only in how the
    generation step is organized, this constitutes an ablation of the multi-agent
    decomposition itself.
\end{itemize}

We additionally relate our measurements to the results published for CodaMosa
\cite{lemieux2023codamosa} and CoverUp \cite{altmayer2025coverup} on this benchmark
family. Neither is treated as a baseline: CoverUp could not be executed to completion
under our local-inference setup, and both are driven by frontier models served through
hosted APIs rather than by the locally served open-weights models used throughout our
experiments. We use the CodaMosa figures to verify that our re-execution of Pynguin is
faithful to the published configuration (Section~5.5), and we report the CoverUp
figures as an end-to-end reference point: the subjects coincide, so the numbers are
comparable as system results, but the difference in underlying model (GPT-4o against
a locally served open-weights model) precludes attributing the gap to architecture.


% ---------------------------------------------------------------------
% [8] INTRODUÇÃO — 1.ª contribuição (substituição integral do bullet)
% PORQUÊ: hoje a contribuição 1 é "o desacoplamento", e é ao desacoplamento
% que o título, o abstract e a conclusão atribuem todos os resultados. Mas
% os dois números mais fortes têm outras causas: os oráculos vêm da política
% de retenção (§3.2.4) e o custo vem do batching ao ficheiro (§5.4). Nenhuma
% experiência isola o desacoplamento sozinho. Reivindicar o PIPELINE — que é
% o que de facto construíste e mediste — é simultaneamente mais honesto e
% mais defensável do que reivindicar uma peça e ser apanhado a atribuir-lhe
% o mérito das outras.
% ---------------------------------------------------------------------
\item A decoupled, multi-agent generation pipeline (\textbf{MARTA}) for Python, whose
effectiveness rests on four mechanisms operating together: the separation of abstract
test planning from concrete syntax formulation, which prevents the context dilution
and instruction forgetting that degrade single-prompt generators; file-level
generation, in which one inference call emits a complete test file for a function;
an inner ReAct loop that repairs the file from execution traces; and a retention
policy that discards what cannot be repaired instead of weakening it. We evaluate the
pipeline as a whole and attribute each measured effect to the mechanism responsible
for it, rather than to the decomposition alone.


% ---------------------------------------------------------------------
% [9] ABSTRACT — última frase (substitui de "We conclude that decoupling..."
%     até ao fim do parágrafo)
% PORQUÊ: "decoupling ... yields robust suites that detect more faults" é a
% atribuição exata que a §5 não sustenta. A conclusão que os dados SUPORTAM
% é sobre o pipeline e sobre a métrica pela qual estes sistemas devem ser
% julgados — e essa continua a ser forte.
% ---------------------------------------------------------------------
We conclude that a pipeline which plans before it writes, generates at file
granularity, and discards what it cannot repair produces suites that detect more
faults per unit of coverage. More broadly, these findings indicate that the value of
such decomposition is revealed by fault detection rather than by structural coverage
alone, a metric a test suite can satisfy without verifying anything at all.


% ---------------------------------------------------------------------
% [10] CONCLUSÃO — 1.º parágrafo (substituição integral)
% PORQUÊ: mesma razão do [9]; e acrescenta a ressalva que desarma a crítica
% em vez de a absorver. Um parágrafo que já reconhece o limite da atribuição
% não pode ser atacado por não o reconhecer.
% ---------------------------------------------------------------------
In this paper, we presented \textbf{MARTA}, an automated code-level robustness and
regression test generation pipeline for Python. By separating abstract planning from
syntax formulation, generating one complete test file per function, repairing it
through a ReAct self-healing loop, and discarding what the loop cannot repair, our
approach translates complex semantic context into executable suites that retain their
full verifying power. Our empirical evaluation shows that this pipeline reduces hollow
oracles by an order of magnitude relative to a single-prompt architecture on the same
context and the same model, attains the highest coverage and mutation score among the
evaluated systems, and generates its suites in half the time. We attribute these
effects to specific mechanisms rather than to the architecture in the abstract: the
oracle gap follows from the retention policy, the cost reduction from file-level
batching, and the coverage gain from the outer planning loop. Isolating the
contribution of the agent decomposition on its own would require an ablation in which
a single agent performs both planning and syntax formulation under otherwise identical
conditions, which we leave to future work.
